% Appendix E: where the time goes in the stock engine, and the serving baseline. Nothing here
% runs the certified head. Byte counts: derived from the checkpoint. Attribution, bandwidth,
% head microbenchmark, host gaps: register entries profile-*. Serving protocol: design
% implemented in bench/; workload, natural lengths, probes, frontend diagnosis, tuning,
% confirmed frontier, output equality, quality: committed (bench-*). Equality against the
% floor: state-pinned, state-perturb. Served certified head: pending (bench-cert). The
% ceilings, levers, P4 and P5 moved to the research notes.
\section{Where the time goes, and the serving baseline}\label{app:stack}\label{sec:stack}\label{sec:time}
This appendix measures the stock engine, the baseline a certified head would have to improve; none of its runs uses the certified head. A component that occupies a fraction $f$ of a step can speed the step up by at most $1/(1-f)$, so the head's share of the time bounds what any head certificate can deliver. Appendices~\ref{app:bytes} and~\ref{app:profile} measure where the time of a plain decode step and of an MTP speculative cycle goes, and Appendix~\ref{app:serving} defines the serving protocol behind the stock frontier of Figure~\ref{fig:frontier}. What the other layers of the stack could deliver, the first lever measurements and the proposals for the recurrent state (P4, P5) are in the research notes.

\subsection{Bytes per step}\label{app:bytes}\label{sec:time-bytes}
Decode on this GPU streams bytes, so the byte counts of Table~\ref{tab:bytes} give a first picture before any trace. A plain decode step streams every decoder weight once plus the whole head, while the embedding lookup reads one entry per request, so the head is 15.1\% of the weight bytes of a plain step. An MTP draft step runs one small layer followed by the same head, so there the head is 84.1\% of the weight bytes. The recurrent state is per request: the configuration stores it in FP32 (\code{mamba\_ssm\_dtype}), so each request holds 50.3~MB of GDN state, and a decode kernel that reads and writes it once per step moves 100.7~MB per request per step, which equals the weight bytes of a step at a batch of 84 requests.

\begin{table}[htbp]
\small\centering
\begin{tabularx}{\linewidth}{@{}P{6.2cm}rY@{}}\toprule
Component & Bytes & Derivation\\\midrule
Tied head (and input embedding) & 1.271 GB & $248{,}320\times2{,}560\times2$ (BF16)\\
Decoder weights other than the embedding & 7.140 GB & sum of tensor sizes in the checkpoint\\
Weights streamed per plain decode step & 8.411 GB & decoder weights plus the head; head share 15.1\%\\
MTP layer & 0.241 GB & sum of tensor sizes; a draft step streams it and the head (head share 84.1\%)\\
Final decoder layer's MLP & 0.142 GB & $3\times2{,}560\times9{,}216\times2$\\
Vision tower (loaded, unused) & 0.667 GB & sum of tensor sizes\\
GDN recurrent state per request & 50.3 MB & $24\times32\times128\times128\times4$ (FP32)\\
Attention KV per context token & 32 KiB & $8\times4\times256\times2\times2$ (BF16 K and V)\\\bottomrule
\end{tabularx}
\caption{Byte counts for Qwen3.5-4B at revision \modelpin, derived from the checkpoint's safetensors headers and \code{config.json}~\cite{qwenconfig}; calculations, not measurements.}
\label{tab:bytes}
\end{table}

\subsection{Measured attribution}\label{app:profile}\label{sec:time-measured}\label{sec:time-method}\label{sec:time-head}\label{sec:time-state}
\paragraph{Method} The attribution uses the production configuration (CUDA graphs for prefill and decode, the overlap scheduler, FlashInfer attention, greedy decoding; Appendix~\ref{app:pins}) with plain decoding at batch sizes 1, 8, 32 and 128 and MTP speculation (three steps, top-$k$ 1, four draft tokens, otherwise default flags, not tuned) at 1, 8 and 32. For each batch size $B$ the client holds $B$ streaming requests with chat-templated essay prompts of 26--36 tokens and a budget large enough that none finishes, and a two-second window is profiled once decode has settled, so the batch is exactly $B$ and no prefill falls inside the window; mean contexts in the plain windows were 987, 862, 688 and 418 tokens. Nsight Systems~\cite{nsight} traces each window with node-level CUDA-graph tracing. Each step is reassembled from the correlation identifiers of its graph replays: a plain step runs from one target-graph replay to the next, and a speculative cycle from one draft-graph replay to the next. Kernels are grouped by name, GEMMs are labelled by their neighbours in the replay, and the labels were checked structurally against the per-step counts the model implies (32 MLP gate/up and down, 48 GDN input and 24 output projections, 8 attention projections of each kind and one head per replay). Where kernels on two streams overlap, each instant is split equally among them, so the categories and the idle time sum exactly to the step~\ev{profile-attr}. Every configuration was also measured with no profiler attached, three windows each, which gives the throughput numbers: plain decoding produced 289, 2,048, 6,298 and 14,388 output tokens per second at $B=1$, 8, 32 and 128 (3.46, 3.91, 5.08 and 8.90~ms per step), with a standard deviation of 0.1--0.4\% of the mean across windows. The collected plain steps are within 2.3\% of the uncollected ones, so their kernel shares are reported directly. The bandwidth reference is measured on this GPU: the best of nine launch configurations of a read kernel sustains 3.83~TB/s over 4~GiB (the nine span 3.59--3.83~TB/s), the same configuration reads exactly the head's 1.27~GB at 3.79~TB/s, and the attribution uses the latter as the peak~\ev{profile-bw}.

\begin{table}[htbp]
\small\centering
\begin{tabular}{@{}lrrrr@{}}\toprule
Share of a plain decode step & $B=1$ & $B=8$ & $B=32$ & $B=128$\\\midrule
Step under the profiler (\textmu s) & 3,541 & 3,982 & 5,128 & 8,886\\
Head chain (head GEMM, FP32 copy of its logits, argmax) & 10.3\% & 9.4\% & 7.8\% & 6.1\%\\
Weight GEMMs (GDN, attention and MLP projections) & 68.0\% & 64.5\% & 54.0\% & 34.5\%\\
GDN recurrent kernel & 3.5\% & 6.8\% & 18.0\% & 41.6\%\\
Other GDN kernels (fused convolution, gated norm, state tracking) & 2.6\% & 2.7\% & 3.3\% & 3.8\%\\
Attention kernels, QK norm and RoPE, gate, KV store & 4.8\% & 5.5\% & 7.2\% & 8.5\%\\
Norms, activation, small kernels and copies & 6.5\% & 7.1\% & 6.7\% & 3.8\%\\
GPU idle, inside and outside graph replays & 4.3\% & 4.0\% & 3.0\% & 1.7\%\\\bottomrule
\end{tabular}
\caption{Where a plain decode step goes, from Nsight Systems windows of 227--572 steps at the stated batch sizes, CUDA graphs and the overlap scheduler on, greedy decoding. Measured~\ev{profile-attr}.}
\label{tab:profile-plain}
\end{table}

\paragraph{Plain decoding} Table~\ref{tab:profile-plain} gives the attribution. The head chain takes 10.3\% of a step at $B=1$ and 6.1\% at $B=128$~\ev{profile-head}, below the head's 15.1\% share of the weight bytes, for two reasons visible in the table. The head GEMM is the most efficient kernel in the step: it streams the 1.27~GB weight at 3.3--3.6~TB/s, while the backbone's weight GEMMs reach 1.6--3.2~TB/s, the split-$K$ projections from 4,096 to 2,560 features being the slowest. And the step contains work that does not scale with the weights: the recurrent state, attention, 65 RMSNorm launches and idle time.

\paragraph{The stock head in isolation} Replaying the exact calls SGLang makes with the real tied weight under CUDA graphs, the head GEMM takes 354.5~\textmu s at $M=1$, 362.0~\textmu s at $M=32$ and 569.2~\textmu s at $M=256$; up to $M=32$ it runs at 93--95\% of the measured read bandwidth, and reading the BF16 weight once at that bandwidth takes 335~\textmu s. cuBLAS selects an \code{nvjet\_sm90\_tst} kernel whose tile depends on $M$. At $M\ge64$ the FP32 copy and the argmax, each re-reading an $M\times248{,}320$ tensor, add 22--50\% on top of the GEMM (the whole chain takes 558.6~\textmu s at $M=128$ and 850.6~\textmu s at $M=256$). Cold L2 changes the GEMM by under 1~\textmu s up to $M=128$ (5~\textmu s at $M=256$), since the weight is 21 times the 60~MB L2~\ev{profile-head}. This chain is the comparator a certified head must beat, and a kernel that reads the same BF16 bytes has at most a few percent to gain at small $M$.

\begin{table}[htbp]
\small\centering
\begin{tabular}{@{}lrrr@{}}\toprule
MTP speculation, three draft steps & $B=1$ & $B=8$ & $B=32$\\\midrule
Cycle under the profiler (\textmu s) & 8,407 & 9,833 & 13,886\\
Cycle without a profiler (ms, mean of 3 windows) & 6.70 & 8.27 & 12.49\\
Output tokens/s without a profiler (plain decoding at the same $B$) & 384 (289) & 2,573 (2,048) & 7,087 (6,298)\\
Accepted length, tokens per cycle (windows without a profiler) & 2.57 & 2.66 & 2.77\\\midrule
Target head chain (verify GEMM over $4B$ positions, FP32 copy, argmax) & 4.4\% & 4.1\% & 3.9\%\\
Draft head chains (three draft GEMMs, FP32 copies, top-1, extend argmax) & 12.9\% & 11.3\% & 8.5\%\\
MTP layer, excluding its head & 4.6\% & 4.4\% & 3.7\%\\
Target weight GEMMs & 29.6\% & 28.4\% & 22.2\%\\
GDN verify kernel (writes four FP32 states per request) & 1.4\% & 5.5\% & 18.3\%\\
GDN state commit (scatter of the accepted state) & 0.5\% & 3.2\% & 8.9\%\\
All other kernels (remainder) & 9.3\% & 10.4\% & 11.0\%\\
GPU idle, inside and outside graph replays & 37.3\% & 32.7\% & 23.5\%\\\bottomrule
\end{tabular}
\caption{Where a speculative cycle goes with the native MTP layer (untuned: three steps, top-$k$ 1, four draft tokens). Shares are of the cycle under the profiler; the idle share is inflated by the profiler, which lengthens this host-bound cycle by 25\%, 19\% and 11\% at $B=1$, 8 and 32 (first two rows). Measured~\ev{profile-attr,profile-head}.}
\label{tab:profile-mtp}
\end{table}

\paragraph{MTP speculation} A cycle replays a draft graph (the MTP layer twice, each time followed by the full head, the FP32 copy and a split-vocabulary top-1), a target-verify graph over four tokens per request, and a draft-extend graph (the MTP layer over the verified tokens and the head once more), with verification and the recurrent-state commit running eagerly between them. Every cycle therefore streams the 1.27~GB head four times. Table~\ref{tab:profile-mtp} gives the attribution. The head takes 17.2\%, 15.3\% and 12.4\% of a cycle at $B=1$, 8 and 32, three quarters of it in the draft heads; because much of the cycle is idle, its share of the time the GPU is busy is larger, 27.5\% at $B=1$~\ev{profile-head}. Most of the idle time is host work in the attention planning of the verify and draft passes~\ev{profile-host}; the research notes trace it and remove part of it. The profiler inflates this gap, chiefly because each graph launch costs about 350~\textmu s under node-level tracing; subtracting the traced GPU-busy time from the cycle measured without a profiler leaves about 1.4, 1.7 and 1.9~ms per cycle idle (21\%, 20\% and 15\%; derived from two measurements). Without a profiler, MTP produced $1.33\times$, $1.26\times$ and $1.13\times$ the plain throughput at $B=1$, 8 and 32 in these fixed-batch windows, which contain no prefill.

\paragraph{What the attribution decides} Other work finds the head a large share of drafting: FR-Spec profiles EAGLE-2 drafting on Llama-3-8B, with a 128k vocabulary, and finds the head projection at 49\% of draft time and the softmax over its output at a further 13\%~\cite{frspec}, a trimmed-vocabulary study finds the head at 64\% of an EAGLE-3-style drafter's FLOPs~\cite{benshoham2026}, and vocabulary sizes have grown with compute budgets~\cite{tao2024vocab}. On this model and GPU the head is a smaller share of the time than of the bytes (H1): removing the head chain entirely is worth at most $1.11\times$ for plain decoding at $B=1$, $1.07\times$ at $B=128$, and $1.21\times$ for the MTP cycle at $B=1$, before any cost of the certificate. These ceilings are Amdahl bounds for removing a component under otherwise unchanged execution, not predicted gains. The other layers of the stack leave more headroom (H6), and the research notes rank them.

\subsection{The serving protocol}\label{app:serving}\label{sec:serving}
Every served speedup is measured against a server, not a kernel. This subsection defines that measurement: what runs, on which prompts, how the latency--throughput frontier is computed, how the baselines are made strong before anything is compared with them, and which quality check a change must pass.

\paragraph{Setup and workload}\label{sec:serving-setup} All runs use the machine and engine of Section~\ref{sec:setup}, serving the BF16 target at its pinned revision with FlashInfer attention~\cite{flashinfer}; the vision encoder, which text requests never reach, uses Triton attention. The client is AIPerf~\cite{aiperf}. Speculative acceptance depends on what the model is writing, so random-token prompts would measure the wrong thing. The benchmark uses a frozen set of 1,857 public prompts in three domains: chat (the first turns of the 80 MT-Bench questions~\cite{mtbench} and English root prompts from OASST1~\cite{oasst1}), code (the 164 HumanEval problems~\cite{humaneval} and MBPP test tasks with their assertions~\cite{mbpp}) and mathematics (GSM8K training questions with a request to box the final answer~\cite{gsm8k}; the GSM8K test split is reserved for the quality check, so no workload prompt is a quality problem). Each source is pinned to a Hugging Face revision and the builder reproduces the files byte for byte. Each domain is cut into disjoint warm-up (43 prompts), tuning (192) and confirmation (384) sets, and every split interleaves the domains so that any prefix of it is balanced, which matters because low-concurrency points use only a prefix. Configurations are chosen on the tuning split and reported on the confirmation split. Templated prompts are short: the confirmation split has a mean of 87.9 prompt tokens (median 78, 99th percentile 283), with chat shorter (mean 47) than code (130) and mathematics (86)~\ev{workload}.

Requests use the model's chat template with thinking enabled (the template's default) and no reasoning parser, so the thinking text is streamed and counted as output. Decoding is greedy. The throughput panel fixes the output at 512 tokens with \code{ignore\_eos}, and every point checks that each response has exactly that length. A fixed length is a faithful sample of the model's behaviour only while the model has not chosen to stop, so a natural-stopping run on 576 tuning prompts measured where greedy outputs end. Only 3.8\% stopped before 512 tokens and 0.7\% had entered a detected repetition loop by then (a 256-token window with fewer than 30\% distinct 8-grams), so almost all of the fixed-length panel is text the model would have produced anyway. The same run shows why greedy decoding is not used for the quality check: 142 of the 576 outputs (25\%; 36\% in mathematics, 32\% in chat, 6\% in code) ran to the 16,384-token limit, a degeneration Qwen's model card warns about for greedy thinking~\cite{qwencard}. The run used the previous workload version, whose mathematics prompts came from GSM8K test, and other processes used up to 24 CPU cores during it, so it is a length measurement, not a timing~\ev{natural-lengths}.

\paragraph{Metrics}\label{sec:serving-metrics} For request $i$ at a sweep point let $n_i$ be the number of generated tokens the server reports in the \path{usage.completion_tokens} field, $s_i$ the time the client sends it, $f_i$ the arrival of its first non-empty streamed chunk and $e_i$ the interval from $s_i$ to its last non-empty chunk. Over the $N$ measured requests the frontier plots
\begin{equation}
x=\frac1N\sum_{i=1}^N\frac{n_i}{e_i},\qquad y=\frac{\sum_in_i}{\max_i(s_i+e_i)-\min_is_i},\label{eq:frontier}
\end{equation}
so $x$ is the mean per-request output rate \emph{including} the time to first token, and $y$ is output tokens per second on the one GPU. Two companions are reported and never substituted for them. The decode rate $x_{\rm dec}=\frac1N\sum_i(n_i-c_i)/(e_i-(f_i-s_i))$, where $c_i$ is the exact token count of the first chunk from per-chunk usage (a speculative step streams several tokens per chunk), excludes queueing and prefill; one token per chunk is never assumed. The steady-state throughput counts only the tokens streamed between the first and the last send, when the client holds its full concurrency, and so excludes the drain at the end of a point. Each point also records time-to-first-token, inter-token and end-to-end latency percentiles, failures, the per-request speculative statistics the server attaches to each stream (verify steps, accepted drafts, and accepted length as output tokens per verify step with the bonus token included), and deltas of the server's counters, among them the fraction of decode passes that replayed a CUDA graph. The scheduler's logged generation rate while the batch is nearly full is recorded as a diagnostic of what the GPU sustains; clients do not see it, and it is never a frontier metric.

\paragraph{Protocol and baselines}\label{app:serving-protocol}\label{sec:serving-protocol} A server arm is a model revision plus a declarative set of launch flags. The harness starts it once, verifies it, and sweeps client concurrency against it in closed loop; the confirmation sweep covers $c\in\{1,2,4,8,16,32,48,64,96,128\}$, each arm over the concurrencies where it competes (\path{bench/campaigns/confirm.sh}). Before trusting a launch it checks that decode CUDA graphs were captured for every batch size up to the server's capacity (the target decode graph, or the target-verify, draft-decode and draft-extend graphs under speculation), that prefill graphs were captured, that the overlap scheduler is running (the scheduler reports it, no fallback is logged, and a stack sample of the scheduler process shows its overlap loop), and that the recurrent-state cache did not reduce the capacity.

Capacity needs care on this model. The attention cache is not the constraint: 32~KiB per token, and the default memory split leaves room for about a million tokens. The recurrent state is. With the prefix cache and the overlap scheduler on, SGLang reserves five FP32 state slots per running request (three for prefix-cache retention and two ping-pong buffers), while decode touches one; the 667 slots the default split gives the state pool therefore cap the server at 133 running requests (code reading, confirmed against the startup log)~\ev{profile-attr}. Speculative verification adds one intermediate state per draft token, and SGLang resets speculative servers to 48 running requests, and logs it, unless the capacity is given; the capacity limits per arm at default settings are \pend{bench-capacity}. Feasibility probes show what launches: plain decoding serves 256 concurrent requests with the prefix cache and 1,024 without it, MTP speculation serves 256 without the prefix cache, and an MTP server with BF16 state at 512 fails at verify-graph capture on a FlashInfer workspace overflow. These probes establish what runs, not performance~\ev{probes}. Every arm is therefore launched at a capacity of 128 running requests, the top of the client sweep (64 for one DFlash arm, below), with the state cache sized for it and the remaining static memory given to the attention cache and the verify states.

Every arm streams every fourth token (\code{stream-interval = 4} in \code{bench/arms.toml}). The server still sends the first token at once and the last at finish, so the times that Equation~\eqref{eq:frontier} uses do not depend on the interval, and the choice takes load off the single-threaded streaming path (below). Each point sends $\max(64,8c)$ measured requests, the first prompts of the confirmation split in order, after one wave of warm-up requests at the same concurrency, and the prefix cache is flushed before every point so that no measured prompt is served from cache. The confirmation sweep runs in three sessions. Each session launches every tuned arm afresh in the same exclusive hold as its matched plain baseline (the FlashInfer plain arm for FlashInfer arms, the Triton one for Triton arms), in reversed order in alternate sessions, so that drift does not align with an arm. Within a session each arm runs once, over its concurrencies in ascending order; the harness alternates the order of concurrencies only between repeats within one launch, which the confirmation sweep does not use. Speedups are ratios within a session; a pair with an invalid point on either side is dropped and counted, never replaced by another run, and points at which the attention cache limited the batch are flagged from the scheduler log. Timed runs hold the GPU exclusively, and because the scheduler, tokenizer manager and client are single-threaded CPU loops, each point also waits until other processes use fewer than two cores and records the foreign load it saw.

The speedup denominator is the strongest existing configuration, not a default. Arm~A is ordinary autoregressive decoding. Arm~B is native MTP speculation, whose draft steps reuse the target's tied head. Arm~C is the public DFlash drafter (Section~\ref{sec:setup}). The search ran on the tuning split~\ev{tuning}: 23 configurations, each launched once and measured at client concurrency 1, 8, 32 and 128 (64 instead of 128 for the two block-16 configurations) with $\max(32,4c)$ requests per point and 512 output tokens. It covered MTP chain depth (two to five steps with the prefix cache, three, five and seven without it), the prefix cache, SGLang's replayed GDN verification and replayed decoding, its batch-adaptive depth, Triton against FlashInfer attention, DFlash blocks of 4, 8 and 16 tokens and FA4 attention for the DFlash drafter. A one-step chain is not in the search, because it failed a launch check that required a draft-decode graph such a chain does not capture; replayed verification was refused at launch for DFlash and with FlashInfer's GDN decode kernel. Draft trees (top-$k$ 2 and 4, which run on this hybrid model, with accepted lengths of 3.45 and 3.56 at concurrency 1 in the feasibility probes~\ev{probes}), the verify attention mode, the prefill chunk size and graph backend were not searched. With one launch per configuration, Table~\ref{tab:tuning} is selection evidence, not a result.

\begin{table}[!t]
\small\centering
\begin{tabular}{@{}llcrrrrc@{}}\toprule
& & & \multicolumn{4}{c}{$y$ (tokens/s) at concurrency $c$} & Accepted\\\cmidrule(lr){4-7}
Arm & Configuration & Prefix cache & 1 & 8 & 32 & 128 & length\\\midrule
Plain & FlashInfer & on & 282 & 1{,}982 & 6{,}087 & 13{,}421 & --\\
& FlashInfer & off & 282 & 2{,}004 & excl. & 13{,}844 & --\\
& Triton & off & 286 & 2{,}018 & 6{,}241 & 13{,}846 & --\\
& FlashInfer, replayed decoding & off & 244 & 1{,}807 & 5{,}909 & 14{,}981 & --\\\addlinespace
MTP & 2 steps & on & 385 & 2{,}288 & 5{,}455 & 9{,}153 & 2.67\\
& 3 steps & on & 461 & 2{,}482 & 5{,}839 & 9{,}593 & 3.28\\
& 4 steps & on & 465 & 2{,}608 & 5{,}761 & 9{,}276 & 3.72\\
& 5 steps & on & 468 & 2{,}575 & 5{,}684 & 8{,}719$^\dagger$ & 4.11\\
& 3 steps & off & 466 & 2{,}534 & 6{,}127 & 10{,}106 & 3.27\\
& 5 steps & off & 481 & 2{,}699 & 6{,}087 & 9{,}456 & 4.14\\
& 7 steps & off & 465 & 2{,}514 & 5{,}690 & 8{,}493 & 4.62\\
& 3 steps, RV & on & 446 & 2{,}646 & 6{,}469 & 11{,}941 & 3.28\\
& 5 steps, RV & on & 462 & 2{,}666 & 6{,}502 & 11{,}186 & 4.11\\
& 3 steps, RV & off & 456 & 2{,}667 & 6{,}925 & 13{,}011 & 3.29\\
& 3 steps, RV, Triton & off & 535 & 3{,}093 & 7{,}103 & 11{,}353 & 3.27\\
& 4 steps, RV, Triton & off & 552 & 3{,}104 & 6{,}958 & 10{,}804 & 3.73\\
& Adaptive depth, RV & off & 474 & 2{,}603 & 5{,}968 & 10{,}010 & 1.97\\\addlinespace
DFlash & Block 4 & off & 498 & 2{,}682 & 6{,}380 & 11{,}477 & 3.28\\
& Block 8 & off & 686 & 3{,}443 & 6{,}727 & 10{,}226 & 4.75\\
& Block 8, FA4 draft attention & off & 769 & excl. & 6{,}983 & 10{,}594 & 4.75\\
& Block 8, Triton & off & 807 & 3{,}613 & 6{,}307 & 8{,}957 & 4.75\\
& Block 16$^\ddagger$ & off & excl. & 3{,}398 & 5{,}386 & 6{,}737 & 5.71\\
& Block 16, Triton$^\ddagger$ & off & 851 & 3{,}639 & 5{,}199 & 6{,}397 & 5.72\\\bottomrule
\end{tabular}
\caption{Configuration search on the tuning split: output throughput $y$ (Equation~\eqref{eq:frontier}) from one launch per configuration with $\max(32,4c)$ requests per point, and the accepted length at concurrency 32. FlashInfer attention unless Triton is named; MTP chains draft one token per step; RV: replayed GDN verification. ``excl.'': excluded because other processes used more than two CPU cores on average, by a sampler that then overcounted, so these may be false positives. $^\dagger$Limited by the attention cache, which without a cap had shrunk to 54K tokens as the state buffers grew: at most 123 requests ran, with retractions. $^\ddagger$Capacity 64; the last throughput column is concurrency 64. Selection evidence, not results~\ev{tuning}.}
\label{tab:tuning}
\end{table}

Three regularities decide the arms. First, speculation pays at low concurrency and loses at high. MTP with three steps delivers 1.64 times plain decoding's $y$ at concurrency 1 (461 against 282 tokens/s); with snapshot verification it falls behind from concurrency 32 (5,839 against 6,087), with replayed verification and no prefix cache it still leads there (6,925) and trails at 128 (13,011 against 13,844), and at concurrency 128 every speculative configuration trails plain decoding. Second, replayed verification, which removes the per-draft-token state snapshots (research notes), recovers 24\% at concurrency 128 for three steps (11,941 against 9,593 tokens/s) at a 3\% cost at concurrency 1, and depth 3 beat depths 4 and 5 at concurrency 128 in every pair the attention cache did not limit, while trailing them by at most 4\% at concurrency 1 and 7\% at 8; in the low-concurrency arm's own configuration (replayed verification, Triton), depth 4 led depth 3 by 3.2\% at concurrency 1 and 0.4\% at 8 and trailed it by 2\% at 32, so one depth serves both MTP arms. Third, Triton attention helps every speculative family at low concurrency and hurts at high concurrency, and plain decoding is indifferent to it. DFlash at block 16 under Triton has the highest throughput of any configuration at concurrency 1 and 8, and at concurrency 1 a per-request rate 3.35 times that of plain decoding under Triton (957 against 286 tokens/s in $x$); among DFlash configurations block 4 leads at concurrency 128 (11,477 tokens/s, 8\% above block 8 with FA4 draft attention).

The chosen arms (\path{bench/arms.toml}) all run without the prefix cache, which this workload does not reuse and whose state slots deeper drafts and replayed verification need. Arm~A is plain decoding under FlashInfer, with a Triton twin as the matched baseline for Triton arms. Arm~B is MTP with three steps and replayed verification, under FlashInfer for high concurrency and under Triton for low. Arm~C is DFlash at block 8 with FA4 draft attention for every concurrency, at block 16 under Triton for low concurrency, capped at 64 running requests because its 16 verify states per request take about 0.8~GB, and at block 4 for high concurrency. Two of these choices change the target's arithmetic: replayed verification computes its verify outputs with a chunked transform on TF32 tensor cores (research notes), and Triton attention for the target replaces the target's attention kernel. Both are exact up to rounding by the comparison below, as is replayed decoding, so every confirmed arm is stock or exact up to rounding. MTP with SGLang's snapshot verification under FlashInfer is kept as a stock reference, and plain decoding with replayed decoding, the fastest configuration at concurrency 128 on the tuning split and the slowest at 1, as a variant.

\paragraph{Output equality}\label{sec:serving-equality} Each arm that changes the target's arithmetic was compared with a matched stock reference on the 320 divergence prompts of Appendix~\ref{app:divergence}, decoded greedily for 256 tokens with the top five log-probabilities at concurrency 1 and without the prefix cache, like every arm: plain levers against stock plain decoding, buffered MTP against stock MTP with three steps, and Triton DFlash at block 16 against stock DFlash at block 16, both DFlash runs at a capacity of 4~\ev{bench-equality}. The class rule was fixed after the first comparisons had been seen. It rests on the divergence classes of Appendix~\ref{app:divergence}, which predate them: an arm is exact up to rounding if every prompt's first divergence is at an exact tie, within one BF16 step, or a near tie with both margins within two steps, and no output differs in length; it is lossy otherwise. Every arm passes, with no event above a near tie. Against its reference, buffered MTP diverges at 3.79 per 1,000 compared tokens, 1.11 times the batch-shape floor of 3.42 (plain decoding at concurrency 1 against 32 on one server with the prefix cache on and at most 16 running requests); its Triton twin at 4.05 (1.19), Triton DFlash at block 16 at 3.82 (1.12), Triton plain decoding at 3.89 (1.14) and replayed plain decoding at 3.40 (0.99). Every ratio's 95\% interval includes 1, and the ratios inform the class rather than decide it. The ratios depend on that floor's cap: with at most 8 running requests and pinned pools the same pair diverges at 0.63 per 1,000 (Appendix~\ref{app:divergence}), while the class rule does not use the floor's rate. The classes above were measured with servers capped at 16 running requests (4 for DFlash), and none of these arms has been compared at a cap of 8; only stock plain decoding and stock MTP were rerun there, and they too diverge only at rounding-level events, an exact tie, one BF16 step or two steps at most~\ev{state-pinned}. Replayed plain decoding was first classed lossy against the divergence study's plain run, which has the prefix cache on: one prompt's first divergence was large. At that position the prefix-cache-on run's own token rests on a logit gap of 0.0625, which stock plain decoding without the prefix cache resolves the other way, and with the prefix cache on a request's log-probabilities can depend on earlier requests (Appendix~\ref{app:divergence}). Against the reference without the prefix cache, replayed plain decoding first differs from it in that prompt at a tie further on, and its rate is the floor's. The two stock plain references, with and without the prefix cache, differ in 96 of 320 prompts at 1.69 per 1,000, all rounding-level.

\paragraph{Speculation against plain decoding} Stock speculation against the same plain reference diverges at 1.15 times the floor's rate both for MTP with three steps (3.95 per 1,000; the ratio's 95\% interval is 0.93 to 1.43) and for DFlash at block 16 (3.92; 0.93 to 1.42). Against the prefix-cache-on run the ratios were 1.29 and 1.28, with intervals above 1, so that apparent excess came at least partly from the reference's prefix-cache setting, not from speculation: at the pin that setting changes the recurrent prefill of prompts longer than one 64-token chunk, from the first output token, in all 6 such prompts of a 40-prompt check of plain decoding at batch 1 and in 192 of the 193 among the 320 for MTP with three steps at batch 1 with pinned pools (Appendix~\ref{app:divergence}). Turning the overlap scheduler off, which at the pin also changes the prefix cache's recurrent-state strategy, changes the same first tokens for MTP. The plain-decoding counterpart on the 320 prompts is \pend{state-mech}. Against this floor, then, an excess of speculation is not detected, though one of up to about 40\% is not excluded. Three limits apply to this check: the floor is a prefix-cache-on pair on one server, and a floor without the prefix cache was not run; the servers' pools were sized from free memory, not pinned, although for MTP with three steps at concurrency 1 without the prefix cache the unpinned run of this comparison and a pinned one at a cap of 8 agree bitwise on all 320 prompts~\ev{state-pinned}; and the DFlash runs used a capacity of 4 against plain decoding's 16. Whether speculation shows an excess depends on the floor, however. With pinned pools and at most 8 running requests, both runs with the prefix cache on, plain decoding at concurrency 1 against 32 diverges at 0.63 per 1,000 while MTP with three steps against plain decoding at concurrency 1 diverges at 3.68, about six times that floor (derived), and every MTP configuration diverges from plain decoding at 3.5 to 4.0 per 1,000, about as often as it diverges from itself between concurrency 1 and 32 (3.1 to 3.5)~\ev{state-pinned}. The speculative rate is about the same in both regimes, and what moves is plain decoding's own batch-shape divergence.

\paragraph{Divergence given perturbation} A perturbation analysis separates how many prompts a change perturbs at all, meaning that their two outputs are not bitwise identical in tokens and top-five log-probabilities, from how often a perturbed output then diverges in tokens~\ev{state-perturb}. Plain decoding's batch shape perturbs 75 of the 320 prompts at a cap of 8, the other 245 being bitwise identical between concurrency 1 and 32, and all 320 at a cap of 16, while speculation perturbs all 320 against plain decoding, at concurrency 1 from the first verify forward. Nearly all perturbed prompts first differ within the first 32 output tokens. Among those, the share whose tokens diverge is 35 of 62 (0.56, 95\% interval 0.44 to 0.68) for plain decoding at a cap of 8, 165 of 318 (0.52, 0.46 to 0.57) at a cap of 16, and 0.51 to 0.57 for every MTP configuration against plain decoding at concurrency 1 and 32; from the first difference on, divergences occur at 3.5 to 4.1 per 1,000 tokens in all of them. The gap between the cap-8 floor (0.63 per 1,000 compared tokens) and the speculative rate (3.5 to 4.0) is therefore mainly the number of prompts perturbed, and on this measure no excess of speculation over either floor is detected. With 62 prompts in that onset range at a cap of 8 the interval is wide, so a modest difference is not excluded, and the intervals assume independent prompts. The conditional rate is not the same for every source of perturbation: over all pairs with at least 30 perturbed prompts the share ranges from 0.29 to 0.60, lowest for MTP under deterministic inference with the verify KV-split patch (28 of 96, 0.29, 95\% interval 0.21 to 0.39, unexplained) and for the FP32 head (130 of 320, which has no exact BF16 ties to flip).

\paragraph{Quality check}\label{sec:serving-quality} Output length and throughput say nothing about whether a change preserved the model. The quality check is the full GSM8K test split (1,319 problems) scored with the GSM8K grader of the \code{sgl-eval} package, version 0.1.2 (a zero-shot prompt, the last boxed answer, symbolic equality) with thinking on, natural stopping at 16,384 tokens, and Qwen's recommended thinking-mode sampling for precise tasks (temperature 0.6, top-$p$ 0.95, top-$k$ 20) under a fixed request seed. Two runs are compared problem by problem: the accuracy difference, an exact McNemar test on the discordant problems, and the share of identical final answers. Sampled accuracy tests equality in distribution for a lossless change and prices a lossy one; it cannot show that greedy outputs are unchanged, which the output comparisons of Appendix~\ref{app:divergence} address. No arm differs detectably from plain decoding~\ev{bench-quality}. Two launches of plain decoding score 89.99\% and 90.30\%, buffered MTP 89.54\%, MTP with snapshot verification 88.93\%, DFlash at block 8 89.69\% and replayed plain decoding 91.05\% (run while it was still classed lossy). Every paired difference against either plain run lies between $-1.36$ and $+1.06$ points, with exact McNemar $p\ge0.18$. A non-significant test excludes nothing by itself. The paired 95\% intervals (Wald intervals on the discordant problems, derived from the committed counts) have lower bounds between $-0.7$ and $-3.3$ points, the lowest for MTP with snapshot verification against the second plain run ($-1.36$ points, interval $-3.27$ to $+0.54$), so the data exclude losses larger than about 3.3 points for every arm and cannot exclude smaller ones. Two sampled runs of one arm agree on the final answer for only 77\% of problems, so the fixed request seed does not make these runs reproducible, and 17 to 19\% of outputs in every arm reach the 16,384-token limit.

\paragraph{Frontiers}\label{sec:serving-frontiers} The baseline frontier is the upper envelope over arms A, B and C, since no arm need be best at every concurrency: speculation saves weight traffic, which matters most at low concurrency, and adds per-request state traffic, which matters most at high concurrency (research notes). Every arm carries an exactness class in \path{bench/arms.toml}: \emph{stock} when it sets only flags known to leave the target's arithmetic as in plain decoding with FlashInfer target attention and stock state handling, \emph{exact up to rounding} by the comparison above, \emph{pending} until compared, and \emph{lossy} otherwise. All nine confirmed arms are stock or exact up to rounding, so the envelope over exact arms equals the envelope over all arms. The frontier over three confirmation sessions, with paired ratios against the matched plain baselines and the run-to-run variance, is committed~\ev{bench-frontier}; only points with three valid sessions rank, so replayed plain decoding, with one session, does not yet \pend{bench-replayssm}. Still pending are the before-and-after frontiers for the certified head in MTP drafting, verification and plain decoding, with output-equality checks at the same batch shape, which decide H4 \pend{bench-cert}, and the exact and lossy stacks of the research notes are pending.

\begin{figure}[htbp]
\centering
% Figure fig:frontier (Section 4.5): the confirmed serving frontier. Its width is the line width up
% to 10.5 cm and its height 0.48 of that. Reads
% evidence/bench/confirm/frontier.csv (per arm and client concurrency: means over sessions
% of x = output tokens per second per request, end to end, and y = output tokens per second
% on the GPU; n = valid sessions) and evidence/bench/confirm/envelope-all.dat (the Pareto
% envelope over the points with at least three valid sessions, grey band). Points with fewer
% than three sessions are hollow; they do not rank on the frontier (evidence/bench/README.md).
% Styles follow style.tex, one colour per arm: plain decoding (arm plain) ink, solid, circles;
% DFlash with 8-token blocks (arm dflash8) violet, solid, diamonds; DFlash with 16-token blocks
% (arm dflash16) reddish purple, dashed, diamonds; MTP (arm mtp) green, dash-dot, triangles,
% filled where three or more sessions are valid and hollow (mark hollow) where only one is, which
% the caption says. Each curve runs from c = 1 (lower right) to c = 128 (upper left); the plain
% curve's ends are labelled, clear of its line (c = 128 above and left of its point, clear of MTP's c = 128
% triangle below it and of the top frame). Error bars are one standard deviation over sessions (x_e2e_std, y_std); every
% one is below 2% of its mean and hidden by its marker, which the caption says. No arm uses the
% certified head. The figure has one panel, so it has no title inside it; the caption names it
% and its status.
% Direct labels, no legend, each placed where its curve is alone (positions checked against
% frontier.csv): plain decoding left of its c = 1-2 run (nothing lies left of x = 228 below
% y = 3,600), opposite the c = 1 label; MTP, two lines, right of its c = 1 point (461, 456), under the DFlash c = 1
% points; 8-token DFlash above and right of its c = 16 point (370, 5,256), where the envelope is
% the upper boundary of every curve; 16-token DFlash under its c = 48 point (130, 5,326), the
% lowest curve left of x = 184, clear of the plain curve at x >= 228.
% Keep only points with at least #1 sessions, or with fewer than #1.
\pgfplotsset{
  only if n at least/.style={x filter/.append code={%
    \ifnum\thisrow{n}<#1\relax\def\pgfmathresult{inf}\fi}},
  only if n below/.style={x filter/.append code={%
    \ifnum\thisrow{n}<#1\relax\else\def\pgfmathresult{inf}\fi}}}
\pgfmathsetmacro{\frontierw}{min(\linewidth,10.5cm)}
\pgfmathsetmacro{\frontierh}{0.48*\frontierw}
\begin{tikzpicture}
\begin{axis}[xmode=log,ymode=log,width=\frontierw pt,height=\frontierh pt,
  xmin=65,xmax=1150,ymin=180,ymax=25000,grid=major,
  xlabel={Per-request rate $x$ (output tokens/s)},
  ylabel={Throughput $y$ (output tokens/s)},
  xtick={100,200,500,1000},xticklabels={100,200,500,1000},
  ytick={300,1000,3000,10000},yticklabels={300,1k,3k,10k},
  % one standard deviation over sessions; keep last in a plot's options (it changes key path)
  sd bars/.style={error bars/.cd,x dir=both,x explicit,y dir=both,y explicit,error bar style={thin}}]
\addplot[line width=4pt,ccontext!25,forget plot,no markers]
  table[x=x,y=y]{../evidence/bench/confirm/envelope-all.dat};
\addplot[arm plain,discard if not={label}{plain-tuned},sd bars]
  table[x=x_e2e_mean,y=y_mean,x error=x_e2e_std,y error=y_std,col sep=comma]{../evidence/bench/confirm/frontier.csv};
\addplot[arm dflash8,discard if not={label}{dflash-tuned},sd bars]
  table[x=x_e2e_mean,y=y_mean,x error=x_e2e_std,y error=y_std,col sep=comma]{../evidence/bench/confirm/frontier.csv};
\addplot[arm dflash16,discard if not={label}{dflash-tuned-b16},sd bars]
  table[x=x_e2e_mean,y=y_mean,x error=x_e2e_std,y error=y_std,col sep=comma]{../evidence/bench/confirm/frontier.csv};
\addplot[arm mtp,discard if not={label}{mtp-tuned},sd bars]
  table[x=x_e2e_mean,y=y_mean,x error=x_e2e_std,y error=y_std,col sep=comma]{../evidence/bench/confirm/frontier.csv};
% the one-session MTP points drawn again over their filled marks, hollow; the caption names them
\addplot[only marks,cmtp,mark triangle,mark hollow,forget plot,discard if not={label}{mtp-tuned},only if n below=3]
  table[x=x_e2e_mean,y=y_mean,col sep=comma]{../evidence/bench/confirm/frontier.csv};
\node[annot,anchor=west] at (axis cs:300,282) {$c{=}1$};
\node[annot,anchor=south east] at (axis cs:106,14800) {$c{=}128$};
% direct labels (positions in the comment above)
% each label in its arm's colour (style.tex), so the two DFlash labels name their lines unambiguously
\node[annot,anchor=east,text=cstock] at (axis cs:255,450) {plain decoding};
\node[annot,anchor=north west,align=left,text=cmtp] at (axis cs:478,640) {MTP, 3 steps,\\FlashInfer};
\node[annot,anchor=south west,text=cdflash8] at (axis cs:390,5700) {DFlash, 8-token blocks};
\node[annot,anchor=north,align=center,text=cdflash16] at (axis cs:135,3800) {DFlash,\\16-token blocks};
\end{axis}
\end{tikzpicture}%

\caption{The confirmed serving frontier of the stock engine for four of the nine tuned arms, from client concurrency 1 (lower right) to 128, or 64 for block 16 (upper left); each point averages its valid sessions (three or four; hollow: one), and one standard deviation over sessions, below 2\% of each mean, is hidden by the markers \ev{bench-frontier}. Grey band: the envelope over all nine arms' points with at least three sessions. All arms are stock or exact up to rounding; none uses the certified head.}
\label{fig:frontier}
\end{figure}

Figure~\ref{fig:frontier} draws four of the nine arms. On the confirmation split, speculation leads plain decoding through client concurrency 32 and trails it from 48 in every drafter family~\ev{bench-frontier}. At concurrency 1, DFlash with 16-token blocks gives 3.06 times the throughput and 3.45 times the per-request rate of its matched plain baseline, both with Triton attention; the best speculative arm, MTP, gives 0.98 times plain at 48 and 0.86 times at 128. Over the 52 points with three or four sessions, the coefficient of variation of throughput has a median of 0.35\% and a maximum of 1.9\%. These stock frontiers fix the baseline a certified verifier must improve~\pend{bench-cert}.

\paragraph{Measurement hazards}\label{app:serving-hazards}\label{sec:serving-hazards} Two effects outside the GPU shaped this protocol. At high concurrency plain decoding can be limited by the streaming path rather than the GPU; speculative arms stream several tokens per chunk and are less exposed, so a frontend bottleneck would inflate their measured advantage. A diagnosis with one plain server at client concurrency 256 on a quiet host measured the effect~\ev{frontend}. With one-token chunks the tokenizer-manager process (HTTP and the OpenAI serving layer) ran at a full core and the client received about 85\% of the server's full-batch decode rate; streaming every four tokens brought that process to 0.35 cores and gave the same throughput as a non-streaming client, about 92\% of the full-batch rate, the rest being the prefill of each new wave and the drain, which $y$ includes by definition. More tokenizer and detokenizer workers made it worse, and more client workers changed nothing. Every row of that diagnosis is a single point; concurrency above 256 on a quiet host is \pend{bench-frontend}. The same diagnosis attributes the low throughput of the earlier feasibility probes at concurrency 256 to 1,024 to CPU contention on the host, not to a frontend limit. That is the second effect: exclusive use of the GPU does not keep CPU-heavy jobs off the host, which is why every point records foreign CPU load and waits for a quiet machine.
